\Needspace{20\baselineskip}
\section{Actual graph labels by fresh vertex naming}
\label{sec:paper-encoding}

This section resolves the representation item that was left open by
Section~\ref{sec:set-presentations}. The repair changes the names of vertices
inside box contents, rather than discarding their binder counts while leaving
the contents literally unchanged. The translated box label is an actual graph.
Its sole input has a name from which its external signature can be read.
Equal graph labels therefore have equal signatures, as required by the paper.

There are no added vertices or edges. All old incidence, wire types, primitive
labels, and ordered boundaries are preserved by explicit bijections, recursively
inside every box. Both representation round trips are proved. Encoding preserves
and reflects recursive isomorphism; encoding and decoding are mutually inverse
on graph classes. The construction accepts every typed hypernet, including the
previous counterexample, and commutes with increasing the depth allowance.

\subsection{Finite sets of actual vertex names}
\termdef{vertex-name}
\termdef{named-graph}
\begin{definition}[Named graph]
Take $\mathbb N^3$ as a common set of vertex names. A named graph stores a finite
list of such names, with vertex carrier the subtype of names occurring in the
list. Thus a name occurring twice in the enumeration is still only one vertex.
The edge carrier is an arbitrary finite carrier. Vertex labels, edge labels,
ordered incidence, and ordered boundary maps satisfy the same conditions as
\texttt{SetNet}. All boundary vertices have ordinary names in this common space.
The operation \texttt{toSet} exposes this subtype as the graph's vertex carrier.
\end{definition}
Its finite enumeration is the list of names equipped with their membership
proofs. Every vertex is in that enumeration by its defining membership proof.
All other fields of the set presentation are inherited unchanged.
\lean{Graph/NamedPresentation}{8}{8}{Graph.VertexName}

\lean{Graph/NamedPresentation}{12}{34}{Graph.NamedNet}

\lean{Graph/NamedPresentation}{41}{64}{Graph.NamedNet.toSet}

\termdef{fresh-naming}
\begin{definition}[Fresh naming]
Given an exhaustive enumeration of a finite carrier, assign each vertex the
position of its first occurrence. For chosen natural numbers $a,b$, give it the
name $(a,b,\mathsf{index}(v))$. The finite name list is the image of the old
enumeration under this function. The map \texttt{naming} sends a vertex to this
name with its membership proof; its inverse chooses a preimage.
\end{definition}
\begin{proof}[Injectivity and inverse laws]
If two first-occurrence indices coincide, looking up the enumeration at that
index returns both vertices, so the vertices coincide. Equality of names implies
equality of their third coordinates, hence equality of vertices. Every member
of the finite name list has a preimage by membership in a mapped list. Choosing
one gives a right inverse; injectivity makes it a left inverse as well.
These facts prove \texttt{index\_injective} and every field of the naming
bijection. Exhaustiveness is essential; duplicate enumeration entries are harmless.
\end{proof}
\lean{Graph/NamedPresentation}{67}{69}{Graph.NamedNet.index}

\lean{Graph/NamedPresentation}{71}{76}{Graph.NamedNet.index_injective}

\lean{Graph/NamedPresentation}{79}{80}{Graph.NamedNet.name}

\lean{Graph/NamedPresentation}{82}{83}{Graph.NamedNet.freshNames}

\lean{Graph/NamedPresentation}{86}{99}{Graph.NamedNet.naming}

\begin{definition}[Freshened set graph]
The operation \texttt{ofSet} changes a set graph's vertex carrier to its fresh
name set. It keeps the edge carrier and edge labels, maps the naming bijection
over every incidence list and boundary vertex, and reads the new vertex labels
through the inverse bijection.
\end{definition}
\begin{proof}[Verification of the graph laws]
Injectivity transfers boundary injectivity, disjointness, and duplicate-free
incidence. Equal named occurrences decode to equal old occurrences, so old
incidence uniqueness identifies their edges. A named vertex belongs to a mapped
incidence list exactly when its inverse belongs to the old list. Combined with
the inverse laws for the boundary maps, this transfers both coverage equivalences.
For typing, composing vertex labels with the naming map cancels the inverse
bijection; the old signature equations then apply.
\end{proof}
\lean{Graph/NamedPresentation}{103}{157}{Graph.NamedNet.ofSet}

\begin{lemma}[Freshening changes only names]
The original set graph and its freshening are label-preservingly isomorphic.
Every freshened input has the chosen first two name coordinates. Starting with
a tagged graph, fresh naming followed by boundary normalization also returns an
ordinary isomorphic graph.
\end{lemma}
\begin{proof}
Use the naming bijection and the identity edge map for the first isomorphism;
all incidence equations hold by construction and labels by the inverse law.
The input-prefix statement is definitional. For the typed round trip, compose
the old tagged-to-set normalization isomorphism with the normalization of the
freshening isomorphism. Literal label equality is reversed where the ordinary
\texttt{Net.Iso} convention requires it.
\end{proof}
\lean{Graph/NamedPresentation}{160}{168}{Graph.NamedNet.ofSet_iso}

\lean{Graph/NamedPresentation}{171}{172}{Graph.NamedNet.input_prefix}

\lean{Graph/NamedPresentation}{174}{184}{Graph.NamedNet.typed_roundtrip}

\subsection{Isomorphism infrastructure for transported boundaries}
\termdef{boundary-list-iso}
\begin{definition}[Boundary-list isomorphism]
\texttt{HRelIso} relates set presentations whose interface indices may have
different presentations. It has bijections of full vertex and edge carriers,
preserves vertex labels and ordered incidence, and relates edge labels by a
supplied relation $R$. It compares the lists enumerated by the boundary maps:
mapping the vertex bijection over each original boundary list gives the new
boundary list. This avoids a premature cast between finite index types.
\end{definition}
\lean{Graph/NamedPresentation}{194}{203}{Graph.SetNet.HRelIso}

\begin{lemma}[Indexed and list formulations]
An indexed related-label isomorphism gives a boundary-list isomorphism. A
literal-label isomorphism is a special case. Conversely, for identical interface
types, a boundary-list isomorphism gives an indexed one.
\end{lemma}
\begin{proof}
For the forward map, mapping a function over a list generated by finite indices
is the list generated by the composite function; use the indexed boundary
equations pointwise. Literal equality supplies the special-case label relation.
For the converse, look up each index in the preserved boundary list. The two
lookups agree, giving exactly the indexed boundary equation. Carrier maps and
incidence fields are unchanged in all three constructions.
\end{proof}
\lean{Graph/NamedPresentation}{209}{217}{Graph.SetNet.HRelIso.ofRelIso}

\lean{Graph/NamedPresentation}{219}{228}{Graph.SetNet.HRelIso.ofIso}
\lean{Graph/NamedPresentation}{322}{334}{Graph.SetNet.HRelIso.same}

\begin{lemma}[Identity, composition, inversion, and relabelling]
Boundary-list isomorphisms have identities for reflexive label relations,
compose for transitive relations, and invert for symmetric relations. More
generally, label relations $R,S$ can compose to $T$ whenever
$R(x,y)\wedge S(y,z)\Rightarrow T(x,z)$. A signature-preserving translation of
labels transports an isomorphism whenever it transports the label relation.
\end{lemma}
\begin{proof}
Identity uses identity carrier maps. Composition composes carrier bijections;
the list equations follow by mapping twice and composing the vertex functions.
The edge-label proof is precisely the specified relational implication.
For inversion, map the inverse vertex function over each original list equation
and cancel the forward map; inverse edge indices cancel by the edge bijection
law. Apply symmetry to the edge-label comparison. Relabelling retains all
carrier and incidence fields and applies its supplied implication to each
edge-label proof. The two composition declarations are the same argument,
first for one relation and then for three potentially different relations.
\end{proof}
\lean{Graph/NamedPresentation}{230}{247}{Graph.SetNet.HRelIso.trans}
\lean{Graph/NamedPresentation}{268}{285}{Graph.SetNet.HRelIso.comp}

\lean{Graph/NamedPresentation}{288}{304}{Graph.SetNet.HRelIso.symm}

\lean{Graph/NamedPresentation}{307}{319}{Graph.SetNet.HRelIso.relabel}
\lean{Graph/NamedPresentation}{374}{382}{Graph.SetNet.HRelIso.refl}

\begin{lemma}[Round trips and independent fresh names]
A named graph is isomorphic to its normalization followed by fresh naming.
Exposing a typed related-label isomorphism gives a full-vertex boundary-list
isomorphism. Independently freshening both sides preserves such an isomorphism.
\end{lemma}
\begin{proof}
The first statement composes the outer set round trip with freshening. For
exposure, map internal vertices by the original bijection and fix every boundary
index; the three port constructors verify labels and boundaries. For independent
freshening, compose the inverse freshening on the left, the given isomorphism,
and the forward freshening on the right. The mixed-relation composition law
uses equality to discharge the two new label comparisons.
\end{proof}
\lean{Graph/NamedPresentation}{255}{258}{Graph.NamedNet.roundtrip}
\lean{Graph/NamedPresentation}{342}{354}{Graph.Net.RelIso.toSet}
\lean{Graph/NamedPresentation}{362}{367}{Graph.NamedNet.map_iso}

\begin{lemma}[Relabelling coherence]
Signature-preserving relabelling of a named graph retains all actual vertex
names. Fresh naming commutes with relabelling. Pointwise equal label translations
induce equal relabelled set graphs.
\end{lemma}
\begin{proof}
Relabelling copies every field except the edge labels and their typing proofs;
the new typing equation composes the old equation with signature preservation.
The fresh name list and naming bijection depend only on the old vertex carrier
and its enumeration, so the commuting equation is definitional. Finally,
function extensionality identifies pointwise equal label maps, after which the
graph records agree by proof irrelevance.
\end{proof}
\lean{Graph/NamedPresentation}{389}{413}{Graph.NamedNet.relabel}
\lean{Graph/NamedPresentation}{420}{423}{Graph.NamedNet.ofSet_relabel}
\lean{Graph/NamedPresentation}{430}{438}{Graph.SetNet.relabel_congr}

\subsection{Boxes labelled solely by actual contents}
\termdef{graph-body-package}
\begin{definition}[Body package]
A body package is a named graph with input type $\star$ and output type
$\star^r$, together with the actual output count $r$. This count describes the
graph's own interface. It is not the bound count or free count of an enclosing
box. Indeed, \texttt{body\_output\_length} proves that it equals the length of
the body's ordered output list, simply by the length of a finite-index list.
\end{definition}
\lean{Graph/PaperEncoding}{13}{13}{Linear.PaperEncoding.Body}
\lean{Graph/PaperEncoding}{358}{359}{Linear.PaperEncoding.body_output_length}

\termdef{actual-content-label}
\begin{definition}[Actual-content label]
For a body package $C$, read numbers $b(C),f(C)$ from the first two coordinates
of the name of its sole input vertex. A content label consists of $C$ and a
proof that $b(C)+f(C)=r$. There are no separately stored binder fields.
Its external signature is $T^{b(C)}(\star)\Rightarrow\star^{f(C)}$.
The base alphabet is the primitive alphabet; a successor alphabet adds these
actual-content labels. A hypernet is a named graph over that alphabet.
\end{definition}
The displayed name format chooses convenient representatives. It does not
restrict the old typed graph classes: arbitrary finite carriers admit the proved
naming bijection, and the encoder below chooses the needed prefix separately
for each box. Wire labels and incidence are untouched by that choice.
\lean{Graph/PaperEncoding}{16}{16}{Linear.PaperEncoding.bound}

\lean{Graph/PaperEncoding}{17}{17}{Linear.PaperEncoding.free}

\lean{Graph/PaperEncoding}{21}{21}{Linear.PaperEncoding.Content}

\lean{Graph/PaperEncoding}{24}{29}{Linear.PaperEncoding.nextAlphabet}

\lean{Graph/PaperEncoding}{31}{33}{Linear.PaperEncoding.alphabet}

\lean{Graph/PaperEncoding}{35}{35}{Linear.PaperEncoding.Hypernet}

\begin{lemma}[Equal contents have equal signatures]
If two content labels have literally equal body packages, their complete external
signatures agree. Forgetting their admissibility proofs adds no label information.
For any graph in the presentation, equal literal labels give equal mapped vertex
labels on both source lists and both target lists.
\end{lemma}
\begin{proof}
Subtype extensionality identifies content labels with equal underlying packages;
substitution proves both signature equalities. Define \texttt{bareLabel} to keep
a primitive or the actual body package and omit the proof. It is injective:
equal primitives are handled by constructor injectivity, mixed constructors
cannot agree, and equal body packages identify their content subtypes.
For two edges, injectivity therefore identifies their original labels.
Rewrite the two incidence-typing equations with that equality. This is the
paper's literal equal-label/equal-signature condition, now checked directly for
graph-valued box labels. It is not a claim that arbitrary isomorphic contents
must give the same enclosing type.
\end{proof}
\lean{Graph/PaperEncoding}{39}{44}{Linear.PaperEncoding.same_content_signature}
\lean{Graph/PaperEncoding}{330}{332}{Linear.PaperEncoding.bareLabel}

\lean{Graph/PaperEncoding}{335}{343}{Linear.PaperEncoding.bareLabel_injective}

\lean{Graph/PaperEncoding}{347}{354}{Linear.PaperEncoding.equal_labels_same_types}

\subsection{Encoding, decoding, and recursive round trips}
\begin{definition}[Signature-preserving alphabet map]
An alphabet map consists of a function on symbols and proofs that it preserves
the complete ordered input and output signatures. These proofs make relabelling
of well-typed graphs possible without changing their incidence.
\end{definition}
\lean{Graph/PaperEncoding}{47}{50}{Linear.PaperEncoding.AlphabetMap}

\begin{definition}[Encoding]
The base alphabet map is the identity. At a successor level a primitive remains
unchanged. An annotated box with counts $b,f$ becomes the actual recursively
translated body, freshly named with prefix $(b,f)$, and with output count $b+f$.
The prefix equation supplies its content-label admissibility proof. Encoding a
whole graph translates its labels and then freshly names its outer vertices
with any chosen prefix $(a,b)$. The prefix of an outer graph is arbitrary;
when the graph is used as box contents, the enclosing box chooses its prefix.
\end{definition}
The input-prefix theorem is definitional. The alphabet-map signature proofs
are also by the primitive/box cases: in the box case, reading the fresh input
name returns exactly the old bound and free counts.
\lean{Graph/PaperEncoding}{53}{63}{Linear.PaperEncoding.encodeMap}

\lean{Graph/PaperEncoding}{67}{69}{Linear.PaperEncoding.encode}
\lean{Graph/PaperEncoding}{90}{92}{Linear.PaperEncoding.encode_prefix}

\begin{definition}[Decoding]
Decoding at depth zero normalizes the named vertex carrier. At a successor it
also translates each content label to an annotated box whose counts are read
from the body's input name. Recursively decode the body, and transport its
output interface along the content proof $b+f=r$. The primitive case remains
unchanged. The resulting alphabet map preserves signatures by construction.
\end{definition}
\lean{Graph/PaperEncoding}{74}{84}{Linear.PaperEncoding.decodeMap}

\lean{Graph/PaperEncoding}{86}{87}{Linear.PaperEncoding.decode}

\begin{theorem}[Typed round trip]
For every typed hypernet $G$, every depth, and every chosen outer prefix,
$G$ is recursively isomorphic to $\mathsf{decode}(\mathsf{encode}(G))$.
Every typed hypernet therefore has an actual-content representative.
\end{theorem}
\begin{proof}
Induct on depth. The base is fresh naming followed by outer normalization.
At a successor use that same outer isomorphism for carriers, incidence, and
vertex labels. A primitive edge label is unchanged. For a box, decoding reads
back exactly the prefix used by encoding, hence recovers both counts. The
body-interface transport is along reflexivity, and induction supplies the
recursive body isomorphism. For representability choose the encoder with outer
prefix $(0,0)$ and this round-trip witness. No body-consistency hypothesis occurs.
\end{proof}
\lean{Graph/PaperEncoding}{96}{112}{Linear.PaperEncoding.roundtrip}
\lean{Graph/PaperEncoding}{363}{365}{Linear.PaperEncoding.representable}

\begin{lemma}[Transport of an invariant property]
Any typed-graph predicate invariant under recursive isomorphism has the same
truth value before and after the decoded encoding.
\end{lemma}
\begin{proof}
Apply the stated invariance hypothesis to the proved typed round-trip
isomorphism. This lemma transports an invariance theorem; it does not assert
invariance for an arbitrary predicate on raw representatives.
\end{proof}
\lean{Graph/PaperEncoding}{116}{119}{Linear.PaperEncoding.roundtrip_property}

\begin{definition}[Recursive actual-content isomorphism]
Primitive labels are compared literally. Two content labels are related when
their derived external counts agree and their actual bodies are related by a
boundary-list isomorphism with recursively related labels. A recursive graph
isomorphism uses this relation for its outer edge labels. Boundary-list
isomorphisms allow the body interface expressions to differ before the
admissibility equations have been used.
\end{definition}
The comparison of counts preserves the enclosing external types. It does not
identify boxes of different external types merely because their contents are
isomorphic as standalone graphs.
\lean{Graph/PaperEncoding}{123}{128}{Linear.PaperEncoding.LabelEquiv}

\lean{Graph/PaperEncoding}{130}{131}{Linear.PaperEncoding.DeepIso}

\begin{theorem}[Named round trip]
Every graph in the actual-content presentation is recursively isomorphic to
its decoding followed by encoding, for any chosen outer prefix.
\end{theorem}
\begin{proof}
At the outer level use normalization followed by fresh naming. Induct on depth
for labels. A primitive returns unchanged. For a content label $C$, the new
body is named with the original derived counts $b(C),f(C)$, so both external
counts agree. Its output interface is transported from $r$ to $b(C)+f(C)$.
A local equality-elimination lemma first handles an arbitrary equation between
these output counts, then invokes induction at the common interface. Thus
neither the cast nor the recursively renamed body is assumed harmless.
\end{proof}
\lean{Graph/PaperEncoding}{135}{160}{Linear.PaperEncoding.named_roundtrip}

\subsection{Faithfulness and equivalence of graph classes}
\begin{theorem}[Isomorphism transport and reflection]
Encoding transports recursive isomorphisms, even with independent outer
prefixes on the two sides. Decoding also transports recursive isomorphisms.
Consequently two encoded typed graphs are recursively isomorphic exactly when
the original typed graphs are recursively isomorphic.
\end{theorem}
\begin{proof}
For encoding, expose the original isomorphism, translate its labels, then apply
independent fresh naming on the two sides. Induct on depth to prove the label
translation respects related labels. Primitive cases are literal equality;
related boxes have equal counts and an induction hypothesis for their bodies.
For decoding, normalize the boundary-list isomorphism at equal outer interfaces.
In its box case, the relation supplies equal bound and free counts. A helper
lemma eliminates these count equalities and both body-interface equations,
leaving the recursive body isomorphism at a common typed interface. Apply
induction there. Reflection now follows by decoding an encoded isomorphism
and composing with the typed round trip on the left and its inverse on the
right; the converse is encoding transport.
\end{proof}
\lean{Graph/PaperEncoding}{164}{183}{Linear.PaperEncoding.encode_iso}

\lean{Graph/PaperEncoding}{187}{223}{Linear.PaperEncoding.decode_iso}

\lean{Graph/PaperEncoding}{226}{231}{Linear.PaperEncoding.encode_iso_iff}

\termdef{representation-class}
\begin{definition}[Representation class]
A representation class is a named actual-content hypernet modulo recursive,
type-preserving graph isomorphism. The equivalence relation uses existence of
such an isomorphism. Encoding and decoding descend to maps between these
classes and the existing typed recursive-isomorphism classes.
\end{definition}
\begin{proof}[Equivalence laws and well-definedness]
Induct on depth for reflexivity, symmetry, and transitivity of label equivalence.
Primitives use the corresponding equality laws. Boxes retain or compose their
derived-count equalities, and their body isomorphisms use identity, inverse,
or composition with the induction hypothesis on labels. These laws supply
the setoid. The quotient maps use the proved encoding and decoding transport
lemmas, so changing a representative cannot change the resulting class.
\end{proof}
\lean{Graph/PaperEncoding}{234}{240}{Linear.PaperEncoding.label_refl}

\lean{Graph/PaperEncoding}{242}{255}{Linear.PaperEncoding.label_symm}

\lean{Graph/PaperEncoding}{257}{275}{Linear.PaperEncoding.label_trans}

\lean{Graph/PaperEncoding}{278}{284}{Linear.PaperEncoding.setoid}

\lean{Graph/PaperEncoding}{286}{286}{Linear.PaperEncoding.Class}

\lean{Graph/PaperEncoding}{289}{292}{Linear.PaperEncoding.encodeClass}

\lean{Graph/PaperEncoding}{295}{298}{Linear.PaperEncoding.decodeClass}

\begin{theorem}[Mutually inverse class maps]
Decoding after encoding is the identity on typed graph classes, and encoding
after decoding is the identity on actual-content graph classes.
\end{theorem}
\begin{proof}
Use quotient induction to choose a representative. The corresponding recursive
round trip supplies the isomorphism witnessing equality of the two quotient
classes. Reverse that equality to obtain the stated inverse law.
\end{proof}
\lean{Graph/PaperEncoding}{301}{305}{Linear.PaperEncoding.decode_encodeClass}

\lean{Graph/PaperEncoding}{308}{311}{Linear.PaperEncoding.encode_decodeClass}

\begin{lemma}[Transport of class predicates and relations]
For every predicate $P$ on typed graph classes, and every binary relation $R$
on those classes, evaluating on decoded encodings has exactly the original
truth value.
\end{lemma}
\begin{proof}
Rewrite each decoded encoding by the first inverse law. For a binary relation
do this at both arguments. No special property of $P$ or $R$ is needed because
they are already defined on classes. A raw context-membership or transition
predicate must first be proved to descend to classes; this result does not
supply an unproved invariance or operational locality law.
\end{proof}
\lean{Graph/PaperEncoding}{315}{318}{Linear.PaperEncoding.class_property}

\lean{Graph/PaperEncoding}{322}{326}{Linear.PaperEncoding.class_relation}

\subsection{Compatibility with larger depth allowances}
\begin{definition}[Promotion of actual contents]
At depth zero promotion moves each primitive to the primitive summand. At a
successor it recursively promotes each actual body while retaining every
vertex name. The input prefix and output count are consequently unchanged,
so the old admissibility proof gives a valid new content label. Relabelling
by this alphabet map promotes a whole graph. Repeated promotion iterates it.
\end{definition}
\lean{Graph/PaperEncodingDepth}{10}{19}{Linear.PaperEncoding.raiseMap}

\lean{Graph/PaperEncodingDepth}{22}{23}{Linear.PaperEncoding.raise}
\lean{Graph/PaperEncodingDepth}{63}{65}{Linear.PaperEncoding.raiseBy}

\begin{theorem}[Exact depth coherence]
Encoding after promotion of a typed graph equals direct promotion of its
encoding. The same equality holds for every finite number of promotions.
\end{theorem}
\begin{proof}
Induct on depth for the one-step assertion. The base is definitional equality.
At a successor isolate the two composed label maps. Fresh naming commutes with
relabelling, so it suffices to prove those maps pointwise equal. Primitive
labels agree immediately. For a box, subtype extensionality removes the proof
field, equality of the output-count package reduces to equality of the actual
bodies, and the induction hypothesis supplies that equality. This proof avoids
expanding the entire naming bijection in each graph field.
Induct on the number of promotions for the iterated statement, applying the
one-step result at the current depth and then the induction hypothesis.
\end{proof}
\lean{Graph/PaperEncodingDepth}{26}{60}{Linear.PaperEncoding.encode_raise}
\lean{Graph/PaperEncodingDepth}{68}{75}{Linear.PaperEncoding.encode_raiseBy}
Together with the common-bound theorem for finite families in
Section~\ref{sec:set-presentations}, this accounts for box contents initially
presented at different finite depth allowances. The equality is for the encoded
graphs themselves, not just their isomorphism classes.

\subsection{Relations on the finite union of depth bounds}
\termdef{global-class}
\begin{definition}[Depth-independent graph class]
A bounded graph is a pair $(n,G)$ where $G$ has depth allowance $n$ and a fixed
ordered outer interface. Define global renaming as the equivalence relation
generated by recursive isomorphism at a fixed bound and the pairs
$(n,G)\sim(n+1,\mathsf{raise}(G))$. Its constructors give these two generators,
reflexivity, symmetry, and transitivity. The global setoid uses exactly these
constructors for its equivalence laws. A depth-independent graph class is a
quotient by this setoid; $\mathsf{global}(G)$ denotes its canonical class.
Operational transitions are not generators of this equivalence.
\end{definition}
\lean{Graph/PaperRelations}{10}{10}{Linear.PaperEncoding.Bounded}

\lean{Graph/PaperRelations}{14}{19}{Linear.PaperEncoding.GlobalRenaming}

\lean{Graph/PaperRelations}{22}{24}{Linear.PaperEncoding.globalSetoid}

\lean{Graph/PaperRelations}{26}{26}{Linear.PaperEncoding.GlobalClass}

\lean{Graph/PaperRelations}{28}{29}{Linear.PaperEncoding.global}

\begin{lemma}[Class invariance]
Recursively isomorphic graphs have equal global classes, and increasing a
graph's depth allowance preserves its class.
\end{lemma}
\begin{proof}
Apply quotient soundness to the isomorphism and promotion generators,
respectively.
\end{proof}
\lean{Graph/PaperRelations}{32}{33}{Linear.PaperEncoding.global_iso}

\lean{Graph/PaperRelations}{36}{37}{Linear.PaperEncoding.global_raise}

\begin{definition}[Interpretation]
Interpret a typed graph by encoding it with outer prefix $(0,0)$ and taking
its global class.
\end{definition}
\lean{Graph/PaperRelations}{40}{41}{Linear.PaperEncoding.interpret}

\begin{proposition}[Interpretation is invariant and surjective]
Interpretation preserves recursive isomorphism and depth promotion. Every
global graph class is the interpretation of a typed graph at some finite bound.
\end{proposition}
\begin{proof}
For isomorphism invariance, encode the given isomorphism and apply class
invariance. For promotion invariance, use the exact encoding--promotion
commuting square, followed by global promotion invariance. For surjectivity,
choose a representative $(n,H)$ by quotient induction and decode $H$.
The named round trip gives an isomorphism from $H$ to its decoded re-encoding;
class invariance and symmetry give the required equality.
\end{proof}
\lean{Graph/PaperRelations}{43}{44}{Linear.PaperEncoding.interpret_iso}

\lean{Graph/PaperRelations}{47}{51}{Linear.PaperEncoding.interpret_raise}

\lean{Graph/PaperRelations}{54}{59}{Linear.PaperEncoding.interpret_surjective}

\termdef{relation-pullback}
\begin{definition}[Relation pullback]
For any relation $R$ on depth-independent classes, define its relation at bound
$n$ by $R_n(G,H)\iff R(\mathsf{interpret}(G),\mathsf{interpret}(H))$.
\end{definition}
\lean{Graph/PaperRelations}{62}{64}{Linear.PaperEncoding.relation}

\begin{proposition}[Relation coherence]
The pulled-back relation is unchanged by independent recursive renaming of
its arguments. It is also unchanged by promoting both arguments once, or by
any common finite number of promotions.
\end{proposition}
\begin{proof}
Unfold the pullback and substitute the equal interpretations supplied by
isomorphism invariance or promotion invariance. For an arbitrary number of
promotions, induct on that number. The zero case is reflexivity; the successor
case composes the induction hypothesis with the one-step equivalence at the
current bound.
\end{proof}
\lean{Graph/PaperRelations}{67}{72}{Linear.PaperEncoding.relation_iso}

\lean{Graph/PaperRelations}{76}{80}{Linear.PaperEncoding.relation_raise}

\lean{Graph/PaperRelations}{83}{90}{Linear.PaperEncoding.relation_raiseBy}

\termdef{global-pretemplate}
\begin{definition}[Global pre-template]
A global pre-template consists of a relation on global classes at each outer
interface, with the requirement that any related typed representatives are
focus-free. This is the defining focus-freedom condition on a pre-template;
it does not supply an operational safety or robustness hypothesis.
Its layered interpretation takes the relation pullback at every depth and
uses the given focus-freedom proof for the existing pre-template interface.
\end{definition}
\lean{Graph/PaperRelations}{93}{96}{Linear.PaperEncoding.GlobalPreTemplate}

\lean{Graph/PaperRelations}{99}{100}{Linear.PaperEncoding.GlobalPreTemplate.toLayered}

\begin{proposition}[Coherent layered interpretation]
The layered interpretation of a global pre-template respects simultaneous
depth promotion and independent recursive renaming.
\end{proposition}
\begin{proof}
Apply relation coherence to its class relation. The layered interpretation
unfolds to precisely that pullback.
\end{proof}
An arbitrary family of unrelated depth-indexed relations is not claimed to
satisfy these properties. This construction supplies a coherent family from
one relation on graph classes; the remaining operational bridges must still
establish their own invariance and transition hypotheses.
\lean{Graph/PaperRelations}{103}{106}{Linear.PaperEncoding.GlobalPreTemplate.coherent}

\lean{Graph/PaperRelations}{109}{112}{Linear.PaperEncoding.GlobalPreTemplate.rename}

\subsection{The repaired counterexample and the remaining theorem}
\begin{proposition}[The two-box obstruction is repaired]
The previous two-box graph admits an actual-content representation without any
restriction on the original graph. The two encoded box labels have different
actual contents, while decoding recovers the old graph up to recursive
isomorphism.
\end{proposition}
\begin{proof}
Encode the unary variable body once with prefix $(0,1)$ and once with prefix
$(1,0)$. If the actual body packages were equal, applying the input-name
bound-coordinate function would give $0=1$. Thus they are distinct graphs.
The first and second edges of the encoded two-box graph carry precisely these
labels by construction. Pair the old negative body-consistency theorem with
the new typed round trip to obtain the repaired-counterexample statement.
The old counterexample remains correct for naive erasure; it is not a
counterexample to freshening.
\end{proof}
\lean{Graph/PaperEncodingExamples}{11}{12}{Linear.PaperEncodingExamples.exposedContent}

\lean{Graph/PaperEncodingExamples}{15}{16}{Linear.PaperEncodingExamples.boundContent}

\lean{Graph/PaperEncodingExamples}{19}{23}{Linear.PaperEncodingExamples.distinct_contents}

\lean{Graph/PaperEncodingExamples}{26}{27}{Linear.PaperEncodingExamples.first_label}

\lean{Graph/PaperEncodingExamples}{29}{30}{Linear.PaperEncodingExamples.second_label}

\lean{Graph/PaperEncodingExamples}{34}{37}{Linear.PaperEncodingExamples.repaired_counterexample}

Further regressions check that a boxed two-hole graph encoded with prefixes
$(2,3)$ and $(7,11)$ gives the same recursive graph class, by applying encoding
transport to the identity isomorphism. The quotient inverse law is instantiated
on that graph. Promoting the repaired two-box graph by two layers commutes with
encoding by the iterated depth theorem. Finally, focus freedom, hole freedom,
and unique-hole contexts are all preserved and reflected by decoded encoding,
using their existing recursive-isomorphism invariance theorems and the new
round-trip isomorphism.
\lean{Graph/PaperEncodingExamples}{40}{44}{Linear.PaperEncodingExamples.prefix_independent}

\lean{Graph/PaperEncodingExamples}{47}{52}{Linear.PaperEncodingExamples.quotient_roundtrip}

\lean{Graph/PaperEncodingExamples}{54}{57}{Linear.PaperEncodingExamples.promoted_repair}

\lean{Graph/PaperEncodingExamples}{61}{66}{Linear.PaperEncodingExamples.context_properties}

\paragraph{Scope of the completed obligation.}
The construction supplies freshly named representatives for the typed UAM
signature of the paper, using its literal graph-label equality in
\S3.1(3). It does not claim a representation theorem for unrelated arbitrary
signatures outside the UAM discipline. Naming is a choice of representative;
no vertices, edges, or operations are added. Box counts are no longer separate
label fields, but can be recovered from names in the actual contents. All typed
graph classes are covered, including graphs rejected by naive binder erasure.

Appendix~\ref{app:remaining-work} removes the representation item and now has
three items. The full-UAM structural results, the concrete specimen interface
with safety/robustness and initial-lifting bridges, and the final application
of Theorem~6 remain. The class-transport lemmas preserve established hypotheses;
they do not prove those outstanding operational hypotheses. The subjective
mathematical-effort estimate is now about 85\%, with a plausible range of
70--91\%, rather than a percentage computed from declaration or to-do counts.
